\section{Related Work and Comparable Projects}

Several studies have used bike-sharing data to compare regression models 
and explore the effect of feature engineering on prediction quality. 
We reviewed five related projects to understand which methods are commonly 
applied and to justify the design choices in our own pipeline.

\begin{table}[htbp]
\caption{Related work and consequences for our project design.}
\centering
\scriptsize
\begin{tabularx}{\columnwidth}{p{0.22\columnwidth} X}
\toprule
\textbf{Source} & \textbf{Relevance and consequence for our design} \\
\midrule
Fanaee-T \& Gama (2013) \cite{fanaee2013} & 
Original study introducing the UCI Bike Sharing Dataset. Confirms 
weather and calendar variables as the standard feature groups for 
this regression task. \\
\addlinespace
UCI Bike Sharing Dataset \cite{uci_bike} & 
Regression benchmark using Capital Bikeshare data from 2011--2012. 
Our project uses a newer dataset (2020--2024) and models demand at 
station-day level instead of city-wide hourly demand. \\
\addlinespace
Kaggle Bike Sharing Demand \cite{kaggle_competition} & 
Many solutions rely on time features, weather variables, and 
nonlinear regression. Supports our pipeline-first approach over 
raw-table modelling. \\
\addlinespace
Capital Bikeshare open data \cite{capital_system_data} & 
Operational data source for station-level demand construction. 
Justifies our station-day aggregation and focus on real operational 
station behaviour. \\
\addlinespace
owenpb \cite{owenpb_github} & 
Kaggle solution comparing Linear, Lasso, Ridge, KNN, Decision Tree, 
Random Forest, and XGBoost. Motivates comparing multiple model 
families rather than reporting only one algorithm. \\
\bottomrule
\end{tabularx}
\end{table}

The reviewed work shows that strong demand forecasts usually depend on 
both model choice and feature construction. This is why our paper gives 
more space to pipeline decisions than to a long list of final plots.